% Avsnitt 2.7, ur lectures/L02/appendix/a_larger_networks.md (A.7) och
% lectures/L02/appendix/c_testbenches.md (C.1-C.6).

\section{Att verifiera din VHDL med en testbänk}
\label{c2:sec:testbenches}

Varje VHDL-modul du skriver till en övning kommer med en \textbf{självkontrollerande testbänk}: en
liten VHDL-fil som driver dina ingångar, kontrollerar dina utgångar och talar om vilket fall som
misslyckades. \Chapref{c1:ch}:s övningar delade redan ut sådana, och \secref{c1:sec:module} gav dig
de tre kommandon som kör en.

Du ombeds inte skriva den testbänk som verifierar en modul du bygger; varenda en av dem delas ut.
Att köra dem räcker för att verifiera varje övning på din egen laptop, utan FPGA-kort. (Ett par
projektövningar från \chapref{c13:ch} och framåt bjuder in till en kort egen testrigg, men
ingenting du bedöms på vilar på en testbänk du skrivit själv.)

Det här avsnittet är bokens introduktion, och varje senare kapitels övningar hänvisar tillbaka
hit. Bilaga~\ref{app:sim} är den fylligare versionen, skriven för grupprojektet; kom tillbaka hit
för introduktionen och gå dit för detaljerna.

\subsection{Det här behöver du}
\label{c2:sub:need}
\textbf{GHDL}, den öppna VHDL-analysatorn och -simulatorn. Kontrollera med \code{ghdl --version};
installera på WSL eller Ubuntu med \code{sudo apt -y install ghdl}.

Varje kommando nedan använder \code{--std=93}, den VHDL-standard boken riktar in sig på.

\subsection{Vad en testbänk är}
\label{c2:sub:what}
En testbänk är själv en VHDL-modul, men en speciell sådan:
\begin{itemize}
  \item Dess entitet har \textbf{inga portar}: ingenting kopplas till omvärlden, eftersom den är
    simuleringens topp.
  \item Dess arkitektur \textbf{instansierar din konstruktion}, ”device under test”, märkt
    \code{dut}, och driver dess ingångar.
  \item Den \textbf{kontrollerar} varje utgång med ett \code{assert}. Om en utgång är fel utlöses
    assertionen med ett meddelande och simuleringen stannar.
\end{itemize}

Så i stället för att du läser av ett vågformsdiagram för att avgöra om utgången är rätt, avgör
testbänken det, och säger bara ifrån när något är fel.

\subsection{Att köra en testbänk: analysera, elaborera, köra}
\label{c2:sub:run}
Tre steg. Skicka med \code{--std=93} till alla tre, och ge \textbf{testbänkens entitetsnamn}, inte
ett filnamn, till de två sista:

\begin{shell}
ghdl -a --std=93 <module>.vhd <module>_tb.vhd        # analyze both files (module first)
ghdl -e --std=93 <module>_tb                         # elaborate the testbench
ghdl -r --std=93 <module>_tb --assert-level=error --stop-time=10ms   # run it
\end{shell}

\begin{itemize}
  \item \code{--assert-level=error} gör att en misslyckad kontroll stoppar körningen och returnerar
    en \textbf{slutkod skild från noll}, ett riktigt godkänt/underkänt som du kan kontrollera med
    \code{echo $?}. Testbänkarna i den här bokens första del reser fel med \code{severity failure},
    vilket stoppar körningen av sig självt, så du får ett riktigt godkänt/underkänt även utan
    flaggan. Ha den på ändå: det är vanan som också fungerar med testbänkar som rapporterar på en
    lägre severity-nivå, och grupprojektets testbänkar är precis sådana (bilaga~\ref{app:sim}).
  \item \code{--stop-time=10ms} sätter ett tak för hur länge simuleringen får köra. En testbänk
    väntar ofta på att din modul ska göra något, med en rad som \code{wait until tx = '0';}. Om din
    modul aldrig gör det blir den väntan aldrig färdig, och utan en stopptid hänger simuleringen:
    ingen utskrift, inget fel, bara en terminal som står och väntar. 10\,ms \emph{simulerad} tid
    är långt mer än någon testbänk här behöver, och tar en bra bit under en sekund i verklig tid.
  \item En godkänd körning skriver ut sin ”all checks passed”-notering och avslutar med 0. Varje
    testbänk i den här boken slutar med den noteringen, så om du inte ser den blev körningen inte
    färdig: den antingen underkände en kontroll eller slog i stopptiden.
\end{itemize}

\subsection{Att kontrollera en övning}
\label{c2:sub:exercise}
Varje VHDL-övning har sin testbänk i sin egen katalog. Det här kapitlets Övning~\ref{c2:ex:xyz}
har till exempel \file{lectures/L02/exercises/xyz_logic/xyz_logic_tb.vhd}. Skriv din modul i
\textbf{samma katalog}, med exakt det namn övningen ber om, och kör sedan de tre kommandona
därifrån:

\begin{shell}
cd lectures/L02/exercises/xyz_logic
ghdl -a --std=93 xyz_logic.vhd xyz_logic_tb.vhd
ghdl -e --std=93 xyz_logic_tb
ghdl -r --std=93 xyz_logic_tb --assert-level=error --stop-time=10ms
\end{shell}

Varje testbänks filhuvud listar precis de här kommandona.

Testbänken kopplar till din modul \textbf{positionellt}, inte via namn, så din entitet måste
deklarera sina portar i \textbf{samma ordning} som övningen listar dem. Namnen är dina att välja,
men använd de angivna ändå, så att din modul och testbänkens felmeddelanden talar om samma
signaler.

En port som deklarerats i fel ordning misslyckas på ett av två sätt, varav inget säger ”fel
ordning”:
\begin{itemize}
  \item om de felplacerade portarna har \textbf{olika} typ eller bredd stannar analysen med
    \code{can't associate "sel" with port "sel"}, med en pekare på testbänkens
    \code{port map}-rad.
  \item om de har \textbf{samma} typ, säg flera \code{std_logic}-portar, analyseras och elaboreras
    den alldeles utmärkt, och faller sedan vid körning som ett vanligt assertion-fel, vilket läser
    sig precis som en logikbugg i din konstruktion.
\end{itemize}

Så när en testbänk faller på allra första fallet den kontrollerar: läs om din portordning innan du
letar igenom din arkitektur. Detsamma gäller \textbf{generics}: en modul som bär på sådana måste
deklarera exakt de som övningen anger, i ordning.

\subsection{När konstruktionen behöver mer än en fil}
\label{c2:sub:multifile}
Vissa övningar bygger en modul av en annan, så \code{ghdl -a}-raden namnger flera filer i stället
för två. Elaborerings- och körstegen är oförändrade: de tar fortfarande bara testbänkens
\textbf{entitetsnamn}.

\begin{shell}
ghdl -a --std=93 <subcomponent>.vhd <module>.vhd <module>_tb.vhd
ghdl -e --std=93 <module>_tb
ghdl -r --std=93 <module>_tb --assert-level=error --stop-time=10ms
\end{shell}

\begin{itemize}
  \item \textbf{Ordningen spelar roll på analysraden}: namnge en modul före allt som instansierar
    den, så att GHDL redan har sett entiteten när det når det \code{entity work.<name>} som
    hänvisar till den. Subkomponenter först, din toppnivå därefter, testbänken sist.
  \item Ingen övningskatalog har någon utdelad subkomponent. Varje modul en konstruktion
    instansierar är en du skrev i en tidigare övning och kopierar in själv; övningen säger vilken,
    och varifrån.
  \item Övningarna som behöver det här är det här kapitlets \code{hex_display} (som instansierar
    din \code{display}), \chapref{c4:ch}:s \code{led_toggle_sync2}, \chapref{c7:ch}:s
    \code{blinker} och \code{walking_led}, samt alla \chapref{c8:ch}:s, där mottagaren namnger fem
    filer och sändaren fyra.
  \item Varje testbänks filhuvud listar det exakta kommandot för sin egen övning.
\end{itemize}

\subsection{Att läsa resultatet}
\label{c2:sub:result}
En misslyckad kontroll ser ut så här:

\begin{console}
xyz_logic_tb.vhd:38:13:@10ns:(assertion failure): xyz_logic: wrong X with a='0' b='0' c='0' d='0', expected '0' but got '1'!
ghdl:error: assertion failed
\end{console}

Läs den i ordning: \code{file:line} är var assertionen står, \code{@time} är när den utlöstes, och
därefter kommer meddelandet. Det är samma form som ett kompileringsfel från GHDL, så det avbildas
på sådant du redan kan. Inget fel betyder att varje fall stämde.

\subsection{Ett komplett exempel}
\label{c2:sub:example}
En kombinatorisk testbänk i den form varje testbänk i den här boken använder, här för
\chapref{c1:ch}:s \code{or_gate}:

\begin{vhdlcode}
library ieee;
use ieee.std_logic_1164.all;
use ieee.numeric_std.all;

entity or_gate_tb is
end entity;

architecture behaviour of or_gate_tb is
signal a, b, x: std_logic;
begin
    dut: entity work.or_gate
        port map(a, b, x);   -- positional: same order as or_gate's port clause.

    SIM_PROCESS: process is
    begin
        for i in 0 to 3 loop
            (a, b) <= std_logic_vector(to_unsigned(i, 2));   -- apply inputs
            wait for 10 ns;                                  -- let them settle
            assert x = (a or b)                              -- check AFTER the wait
                report "or_gate: wrong output with a=" & std_logic'image(a)
                     & " b=" & std_logic'image(b)
                     & ", expected " & std_logic'image(a or b)
                     & " but got " & std_logic'image(x) & "!"
                severity failure;
        end loop;
        report "or_gate: all checks passed!" severity note;
        wait;                                                -- end the simulation
    end process;
end architecture;
\end{vhdlcode}

Två regler håller en kombinatorisk testbänk korrekt:
\begin{itemize}
  \item \textbf{Lägg på, vänta, kontrollera, i den ordningen.} En signaltilldelning
    \emph{schemalägger} bara ett värde, så ingångarna och utgången uppdateras inte förrän vid
    \code{wait}. Att kontrollera före den testar föregående varvs värden.
  \item \textbf{Avsluta med \code{wait;}} så att stimulusprocessen suspenderas och simuleringen
    stannar av sig själv.
\end{itemize}

Klockade konstruktioner behöver lite mer, en klockgenererande process och ett sätt att stoppa den,
vilket är precis det maskineri de utdelade testbänkarna redan innehåller. Till övningarna skriver
du bara modulen.

\paragraph{Kör en nu}
Mot det här kapitlets första genomarbetade exempel:

\begin{shell}
cd lectures/L02/gate_network
ghdl -a --std=93 gate_network.vhd gate_network_tb.vhd
ghdl -e --std=93 gate_network_tb
ghdl -r --std=93 gate_network_tb --assert-level=error --stop-time=10ms
\end{shell}
